% Source: https://artofproblemsolving.com/community/c3h2252795_how_can_an_infinite_series_have_a_sum
% Stable identity: aops-topic-2252795-q01
Let $a\ne0$ and $r\ne1$, and let $S_n=a+ar+ar^2+\cdots+ar^{n-1}$.
\begin{enumerate}
\item[(a)] Show that $S_n=\dfrac{a(1-r^n)}{1-r}$.
\item[(b)] Suppose $|r|<1$. Find $\dfrac{a}{1-r}-S_n$ and explain why $S_n$ gets as close as we like to $\dfrac{a}{1-r}$ as $n$ grows. (This is what is meant by saying that the infinite series $a+ar+ar^2+\cdots$ has sum $\dfrac{a}{1-r}$.)
\item[(c)] For $\dfrac{1}{2}+\dfrac{1}{4}+\dfrac{1}{8}+\cdots$, find the sum and the least $n$ for which the sum of the first $n$ terms is within $0.001$ of it.
\item[(d)] Explain why the series $1+1+1+\cdots$ and $1-1+1-1+\cdots$ ($r=1$ and $r=-1$) have no sum in this sense.
\end{enumerate}
